\section{Results and measured tradeoffs}
\begin{table}[ht]
\centering\small
\begin{tabular}{lrrrr}
\toprule
Direct quartic over & Variables & Rows & Terms & Height\\
\midrule
$\Z$, sparse &215&152&14,578&16,866\\
$\Z$, low height &215&152&34,700&4,512\\
$\N$, sparse &202&154&9,353&1,854,167\\
$\N$, low height &202&154&16,927&416,488\\
$\N$ &203&155&6,524&1,979,888\\
$\N$ &204&156&6,503&1,981,520\\
\bottomrule
\end{tabular}
\caption{Complete degree-four equations. Height is the largest absolute
coefficient after collecting the entire sum of squared residuals.}
\end{table}
The new sparse integer equation improves the incoming $216$-variable,
$153$-row equation with $15{,}158$ terms and height $334{,}952$ in all four
displayed measures. Against our preceding sparse $217$-variable equation,
it saves two variables, two rows and $37$ terms at unchanged height.
The low-height companion trades larger support for height $4{,}512$.
Larger-variable incoming equations with still lower heights remain genuine
tradeoffs; no global coefficient minimum is claimed.

The sparse $202$-natural-variable equation improves both the $9{,}818$
terms and height $239{,}204{,}765$ of its matched predecessor. The lower-height
profile lowers the prior best displayed height $31{,}145{,}294$ to $416{,}488$
at greater support. These changes preserve the complete natural zero set
through explicitly proved coordinate and residual maps.

A separate natural companion has $200$ variables, $152$ residuals,
exact degree $8$, $10{,}392$ terms and height $13{,}654{,}244$.
It eliminates two auxiliary coordinates at a degree cost and is not a
$200$-variable quartic.

\begin{table}[ht]
\centering\small
\begin{tabular}{lrrrl}
\toprule
Domain & Variables & Degree & Operations & Matched comparison\\
\midrule
$\Z_{>0}$ &10&206&90&Late source $214/87$\\
$\Z_{>0}$ &10&210&89&Late source $214/87$\\
$\Z$ &14&85&110&Earlier $86/112$\\
$\Z$ &14&87&108&Earlier $88/110$\\
\bottomrule
\end{tabular}
\caption{Integer-leaf circuit measurements. Comparison pairs are
degree/operations at the same variable count.}
\end{table}
The ten-positive-variable equations are new operation/degree tradeoffs:
they lower degree but add operations relative to $214/87$. The two
fourteen-integer equations improve both degree and operation count over
their matched parents. The incoming $14/82/113$ point remains a different
tradeoff, as do lower-variable equations of higher degree.

The complete Start-marker4177 specialization, credited to its incoming
source, uses $98{,}219$ raw bytes across four coherent coefficient ports;
the largest has $196{,}790$ bits. This is not a new encoding invented here.
The new carrier front is separately bound to its complete native proof
and finite controller certificate.
